% critique.tex — f168v, vv. 509–550 : commentaire critique
% Articles textuels tirés du commentaire de l'auteur (latex/Aristophane-Thesmo Harmattan 2023 vv 1-641.tex),
% texte et notes conservés ; répartition détaillée dans repartition.md.
\cc{509}{τέξειν.}{Cf. sur le mot τέξειν la note 14 de Willems à lire dans les annexes.}
\cc{531–2}{ἀλλ᾽ οὐ γάρ ἐστι τῶν ἀναισχύντων φύσει γυναικῶν / οὐδὲν κάκιον εἰς ἅπαντα πλὴν ἄρ᾽ εἰ γυναῖκες.}{Pour voir ce que Willems pense de la « manie des conjectures » visant à retoucher le texte d'Aristophane, on lira avec un plaisir non dissimulé la note entière n° 15, reproduite à la fin.}
\cc{533}{μὰ τὴν Ἄγραυλον.}{Aglaure (correction d'Agraule confirmée par l'épigraphie).}
